% =====================================================================
%  1. INTRODUCCIÓN
%  Responsable: TODOS
%  ESTADO: REDACTADO (1 de agosto de 2026) — ~950 palabras
%  Se redactó AL FINAL, con el informe ya cerrado, para que el último
%  párrafo describa la estructura real y no la prevista.
% =====================================================================

\section{Introducción}
\label{sec:introduccion}

La minería es una actividad de alto impacto socioambiental por definición y no por
defecto de gestión. Transforma de manera irreversible el relieve, moviliza
volúmenes de material que superan en órdenes de magnitud a los del mineral
aprovechado, consume y puede alterar la calidad del agua de las cuencas donde se
emplaza, y se implanta sobre territorios que rara vez la solicitaron. Su
particularidad frente a otras actividades industriales reside en que no puede
elegir su localización: se instala donde está el yacimiento, y con frecuencia el
yacimiento está bajo el territorio de poblaciones rurales o indígenas cuya
economía depende directamente de la tierra y del agua que la operación va a
utilizar. De esa condición estructural nace la doble dimensión —ambiental y
social— que da nombre a este informe y que ningún instrumento de control puede
abordar por separado sin quedar ciego ante la mitad del problema.

Frente a ese impacto, la auditoría opera como herramienta de verificación. No
previene, como el estudio de impacto ambiental, ni sanciona, como la fiscalización
estatal: comprueba, mediante un proceso sistemático, documentado e independiente,
si la conducta efectiva de la organización se ajusta a criterios previamente
establecidos, y traduce las desviaciones detectadas en hallazgos y recomendaciones
\parencite{iso2018iso19011}. Su valor depende por completo de tres condiciones
que este trabajo examina con detenimiento: que los criterios sean exigentes y
explícitos, que quien audita sea efectivamente independiente de quien es auditado,
y que exista alguien con capacidad y mandato para verificar que lo recomendado se
cumplió. Cuando alguna de las tres falla, el resultado puede ser un documento
técnicamente correcto y prácticamente inocuo.

El caso seleccionado para este análisis es la Mina Marlin, operada por Montana
Exploradora de Guatemala, S.\,A., filial de Goldcorp Inc., en el departamento de
San Marcos, Guatemala, entre 2005 y 2017. La elección responde a cuatro criterios.

El primero es la existencia de un documento de auditoría público, independiente y
exhaustivo: la \hria, elaborada por On Common Ground Consultants Inc. y publicada
en mayo de 2010, con más de doscientas páginas y versión bilingüe
\parencite{oncommonground2010hria}. El segundo, y decisivo, es que ese documento
es \textbf{ambiental y social a la vez}: evaluó siete ejes —consulta, ambiente,
trabajo, adquisición de tierras, inversión económica y social, seguridad y acceso
a remediación— aplicando estándares internacionales de derechos humanos, que es
exactamente el tipo de auditoría socioambiental que este trabajo debe analizar. El
tercero es la disponibilidad de datos técnicos verificables en fuentes primarias:
el informe técnico NI 43-101 presentado ante la Comisión de Bolsa y Valores de los
Estados Unidos \parencite{sec2011ni43101} y la documentación de la Corporación
Financiera Internacional \parencite{ifcsfmarlin}. El cuarto es que el caso tiene
\textbf{ciclo completo}: operación, conflicto, auditoría, plan de acción, cierre y
post-cierre. Ello permite evaluar la mejora continua sobre un expediente cerrado y
no sobre una operación en curso, y comprobar si las recomendaciones formuladas en
2010 llegaron efectivamente a cumplirse antes del cierre de 2017.

Se consideraron y descartaron otros casos. Brumadinho, Mount Polley y Samarco
disponen de investigaciones técnicas de calidad superior en materia geotécnica,
pero sus paneles de expertos son auditorías esencialmente técnicas: excelentes
para el eje ambiental y débiles para el social. Espinar, en el Perú, presenta un
componente social y de diálogo muy fuerte, pero su documento ancla es un monitoreo
sanitario ambiental participativo y no una auditoría en sentido estricto. Marlin
resultó ser el único caso en que el documento auditado combina de forma nativa
ambas dimensiones.

La pertinencia del caso para un lector peruano merece explicitarse, pese a
tratarse de una operación guatemalteca. La conflictividad socioambiental minera en
el Perú responde a una estructura semejante —operaciones de gran escala sobre
territorios de comunidades campesinas e indígenas, disputas sobre calidad y
disponibilidad de agua, y déficits de consulta previa—, de modo que el caso
funciona como término de comparación. Más aún: permite contrastar dos modelos de
control de naturaleza distinta, la auditoría voluntaria financiada por la empresa
auditada frente a la fiscalización estatal con potestad sancionadora que ejerce el
Organismo de Evaluación y Fiscalización Ambiental
\parencite{congreso2009ley29325}. Ese contraste, desarrollado en la sección
\ref{sec:critico}, constituye uno de los aportes centrales del trabajo.

El informe se organiza en diez secciones. Tras estos preliminares y el enunciado
de los objetivos, la sección \ref{sec:marco} construye el marco teórico y
normativo: delimita la auditoría ambiental frente a la fiscalización y a la
evaluación de impacto ambiental, expone las normas ISO 14001:2015 e
ISO 19011:2018 y los estándares internacionales de desempeño social, y presenta el
marco peruano de referencia y el ciclo PHVA. La sección \ref{sec:unidad} describe
la unidad minera: titularidad, ubicación, entorno social, geología, método de
explotación, procesamiento, producción e infraestructura crítica. La sección
\ref{sec:auditoria} reconstruye el proceso de auditoría —origen del encargo,
organismo auditor, equipo, alcance, criterios y metodología— y lo contrasta con
las fases que establece la ISO 19011:2018. La sección \ref{sec:hallazgos}
sistematiza los hallazgos en una matriz de no conformidades y desarrolla los
cuatro de mayor gravedad. La sección \ref{sec:critico} evalúa críticamente el
proceso: independencia, limitaciones metodológicas, conflicto de evidencia sobre
la contaminación, eficacia real y contraste con el régimen peruano. La sección
\ref{sec:plan} formula la propuesta de plan de mejora continua bajo el ciclo PHVA.
Las secciones \ref{sec:conclusiones} y \ref{sec:recomendaciones} recogen,
respectivamente, las conclusiones referidas a cada objetivo específico y las
recomendaciones dirigidas a la empresa, al Estado, a las comunidades y a la
comunidad académica. Cierran el documento las referencias y cinco anexos.
